% Seção 5 — Metodologia. Revisão do Estágio 4 (2026-11-15), após a decisão
% "revisão maior" do painel (docs/ars/pipeline-rodada20/stage3/). Decisões do
% pesquisador: reparar o desenho de calendário; diagnóstico descritivo, sem
% poder de decisão; pré-período de referência fixado (2007–2011); peso ao
% nascer como única hipótese confirmatória. NENHUM número de estimativa aqui.
% Revisão da rodada 21 (2026-11-16), após a decisão do Estágio 3: parâmetro
% como contraste, hipótese de retorno dos voos, regra de decisão completada
% no Apêndice F. Identificação e estimador principal não mudaram.

\chapter{Metodologia}
\label{cap:metodologia}

\section{Parâmetro-alvo e cadeia causal}

O parâmetro de interesse é o efeito da lei, em intenção de tratar, sobre a
diferença de peso ao nascer entre as coortes cujo início de gestação cai na
estação de pulverização e as expostas a ela apenas no fim da gestação, nos
municípios cearenses com aplicação aérea documentada no período de referência.
Com $\Delta_{mt} = Y_{mt}(1) - Y_{mt}(0)$ a diferença entre os desfechos com e
sem a proibição, no pós-lei,
\begin{equation}
  \tau = E\left[\,\Delta_{mt} \mid A_m = 1,\ H_t = 1\,\right]
       - E\left[\,\Delta_{mt} \mid A_m = 1,\ R_t = 1\,\right],
\end{equation}
em que $A_m$ indica município com aplicação aérea documentada entre 2007 e
2011, $H_t$, coorte com exposição alta à estação no início da gestação, e
$R_t$, coorte de referência, exposta apenas no terceiro trimestre; as médias
são por nascimento. O contraste só limita por baixo o efeito em nível sob uma
hipótese mantida, não testável: a de que a exposição tardia não beneficia o
feto. O tratamento é a lei, e não o voo: o que se estima inclui o
descumprimento e a substituição pela aplicação terrestre. A referência de
$A_m$ é anterior à seca de 2012--2017 \cite{rebello2018seca}, em que os voos
podem ter sido suspensos (Seção~\ref{sec:onde-aviao}); definir o tratado pelo
que ocorria às vésperas da lei poderia deixar o grupo vazio. Essa escolha
exige uma hipótese de identificação adicional, não verificada: sem a lei, os
voos teriam voltado ao padrão da referência depois da seca. Sob ela, $\tau$ é
o efeito de a lei ter impedido os voos onde eles ocorreriam. Os relatórios de
2008 a 2024 e o relato de voo até 2018 são a evidência que a examina; se os
voos cessaram antes da lei por outra causa, um coeficiente positivo não
distingue o efeito da lei do fim da prática (Quadro~\ref{qua:regra}, linha 5).

A cadeia que liga a lei ao desfecho tem elos que podem falhar separadamente: a
lei precisa ter sido cumprida; precisava haver voos a suspender ou a impedir
de voltar; a ausência de voos precisa reduzir a exposição das gestantes, e não
apenas trocá-la pela da aplicação terrestre; e a menor exposição precisa
alterar o desfecho. Um efeito nulo no desfecho final não distingue falha de implementação de
falha do mecanismo, e por isso o desenho mede os elos intermediários.

\section{Desenho principal: diferença tripla de calendário}
\label{sec:calendario}

O Apêndice~\ref{ap:dag} traz o diagrama causal do desenho e liga cada
elemento dele à hipótese de identificação correspondente. Dentro de cada município, comparam-se as coortes de
exposição alta à estação no início da gestação com as expostas apenas no fim;
essa diferença, o hiato sazonal, é comparada entre períodos e entre os
municípios com aplicação aérea e os de comparação.

Seja $P_t \in [0,1]$ a fração dos dois primeiros trimestres de gestação da
coorte $t$ que cai na estação de pulverização (convenções no
Apêndice~\ref{ap:espec}). As coortes são classificadas em
três grupos, fixados neste projeto: exposição alta, $H_t = 1$, se $P_t$ é de
metade ou mais; intermediária, $I_t = 1$, se $P_t$ é positivo e menor que a
metade; e nula no início da gestação, a referência, se $P_t = 0$. Sejam ainda
$Seca_t$ o indicador das coortes nascidas de 2013 a 2018, cujo início de
gestação cai, em geral, nos anos de seca, $Pos_t$ o das
nascidas a partir de 2020 e $C_{m,q(t)}$ a chuva, no município, da quadra
chuvosa a que o início da gestação foi exposto. Estima-se, com ponderação pelo número de nascimentos,
\begin{equation}
\begin{split}
  Y_{mt} ={}& \alpha_{m,s(t)} + \mu_{m,a(t)} + \gamma_t
   + \sum_{k \in \{I,H\}} \Bigl[\theta_k\,(k_t \times C_{m,q(t)})
   + \beta^k_S\,(A_m \times k_t \times Seca_t) \\
  & + \beta^k\,(A_m \times k_t \times Pos_t)\Bigr] + \varepsilon_{mt},
\end{split}
\label{eq:ddd}
\end{equation}
em que $\alpha_{m,s(t)}$ são efeitos fixos de município por mês do ano,
$\mu_{m,a(t)}$ de município por ano de nascimento e $\gamma_t$ de coorte. Os
primeiros absorvem a sazonalidade de base de cada município; os segundos, os
choques comuns a todas as coortes de um município num ano, entre eles os
programas estaduais de saúde materna contemporâneos à lei; as coortes de
transição, nascidas em 2012 e em 2019, são excluídas. Os termos de
ordem inferior são absorvidos pelos efeitos fixos. O coeficiente de
interesse é $\beta \equiv \beta^H$: quanto a diferença de peso ao nascer entre
as coortes de exposição alta e as de referência mudou, do período de
referência para o pós-lei, nos municípios com aplicação aérea em relação aos
de comparação. Sob a hipótese principal, $\beta > 0$. O coeficiente
$\beta_S \equiv \beta^H_S$ mede a mesma variação no período de seca e é
descritivo.

A exposição entra em grupos, e não como variável contínua, porque o
coeficiente de uma exposição contínua mistura efeitos de níveis diferentes,
por analogia com \citeonline{callaway2024continuous}
(Capítulo~\ref{cap:referencial}). Com uma
estação de quatro meses, toda gestação a termo a atravessa em algum trimestre:
as coortes de referência têm apenas o terceiro na estação, e por isso o
parâmetro é um contraste, que subestima o efeito em nível se a exposição
tardia também causar dano. A janela gestacional é fixa: a duração observada da
gestação é desfecho, e usá-la para definir a exposição condicionaria em
variável afetada pelo tratamento.

\textbf{Período de referência.} São as coortes nascidas de 2007 a 2011, anos
em que há relatos de voos na Chapada (Seção~\ref{sec:onde-aviao}) e que
antecedem a seca. A escolha foi registrada em 15 de novembro de 2026, depois
da rodada exploratória de setembro e antes da de novembro, que a tomou como
dada (Apêndice~\ref{ap:regra}). As coortes de 2005 e 2006 entram apenas
no estudo de evento. Três limites são conhecidos: os registros pedidos começam
em 2008; em 2011 houve projeto de lei estadual e ação civil pública pela
proibição; e Limoeiro do Norte teve proibição municipal de novembro de 2009 a
maio de 2010. Reportam-se a referência restrita a 2008--2010 e a exclusão das
coortes de Limoeiro afetadas (Apêndice~\ref{ap:espec}).

\textbf{Grupo de comparação.} Os municípios de comparação são os produtores de
banana em que a cultura é predominantemente familiar, segundo o Censo
Agropecuário de 2017 \cite{ibge2017censoagro}, com limiar a fixar no plano de
análise. Supõe-se que não tenham aplicação aérea, o que se verifica nos
registros e no Censo de 2006, que a registra também fora dos polos irrigados;
município com qualquer registro de aeronave é excluído. Em nível, a comparação
não seria válida: agricultura familiar e empresarial diferem em renda, acesso
a programas sociais e exposição à aplicação costal. Dentro da diferença
tripla exige-se menos: que, na ausência da lei, o hiato sazonal tivesse
evoluído do mesmo modo nos dois grupos desde o período de referência
\cite{olden2022triple}. A hipótese é examinada por um estudo de evento,
com um coeficiente por ano e 2011 como base, que fora da referência cobre só
2005 e 2006, com pouco poder, e por um tratamento falso, que compara os
municípios de banana familiar aos sem banana e examina o grupo de comparação,
não o contrafactual dos tratados. Os dois são relatados junto a $\beta$, e não
como portão: um resultado desfavorável não suspende o teste, e $\beta$ é
sempre lido com intervalo de sensibilidade a desvios das tendências paralelas
do tamanho observado antes da lei
\cite{rambachan2023credible}. Como a identificação pode depender de características
dos municípios, reporta-se, se o número de tratados permitir, um estimador
duplamente robusto com covariáveis anteriores à lei \cite{ortiz2025triple}.

\textbf{Chuva.} A chuva entra de dois modos. Como a precipitação durante a
gestação afeta o peso ao nascer no semiárido por vias que não passam pelo
avião \cite{rocha2015water}, a chuva da quadra chuvosa, interagida com os indicadores de exposição, é controle nos dois grupos, sob a hipótese adicional, declarada, de que esse efeito
direto seja comum a municípios de lavoura irrigada e de sequeiro. É ela que
separa a lei do retorno das chuvas; por isso a conclusão confirmatória exige
que a especificação principal e a que deixa o efeito diferir por grupo
concordem em sinal e em ordem de grandeza (Apêndice~\ref{ap:regra}), e, se
divergirem, $\beta$ é reportado como resultado não separado do retorno das
chuvas. E, como a doença que motiva o
voo depende de umidade, examina-se, em caráter exploratório, se $\beta$ é maior
nos anos chuvosos. As especificações estão no Apêndice~\ref{ap:espec}.

\textbf{Diagnóstico descritivo do pré-período.}\label{sec:diagnostico} Antes
da análise confirmatória, descreve-se o hiato sazonal nos dois grupos no
período de referência e no período de seca. O diagnóstico responde se a
exposição à estação aparece nos dados municipais e se a seca a alterou. Ele é
descritivo: não escolhe o período de referência e não autoriza nem impede a
análise confirmatória (Seção~\ref{sec:regra}; a razão está no
Apêndice~\ref{ap:ameacas}). A leitura tem um limite: os dois períodos diferem
pela própria seca, que altera a sazonalidade do peso ao nascer por outras
vias, de modo que o fechamento do hiato nas coortes de 2013 a 2018 é
compatível com a suspensão dos voos, mas não a demonstra.

\section{Medição do método}
\label{sec:medicao}

A classificação $A_m$ e o calendário de pulverização são medidos, e não
supostos. A fonte principal são os relatórios mensais que as empresas de
aviação agrícola entregam ao Ministério da Agricultura \cite{mapa2008in2}, com
município, cultura, área e produto; pedem-se os de 2008 a 2024, para o Ceará e
o Rio Grande do Norte, incluídos os operadores sediados em outros estados, e a
relação de operadores registrados, com a qual se afere a completude da série.
Os anos posteriores à lei medem o cumprimento e, desde 2021, o uso de drones
\cite{mapa2021portaria298}. Não se sabe se os relatórios anteriores ao envio
eletrônico foram preservados, e por isso a entrega da medida é condicional; o
silêncio de uma série incompleta não é lido como ausência de voo (regra de
classificação e critério de completude no Apêndice~\ref{ap:regra}).
Complementam a medida o cadastro estadual de prestadores, os autos de
infração, o receituário agronômico e os autos da ação civil pública de 2011
(Apêndice~\ref{ap:fontes}). O Censo Agropecuário de 2006
\cite{ibge2006censoagro} fornece a classificação provisória; por não
distinguir cultura e ser anterior à seca, pode apontar para municípios
diferentes daqueles em que se voava em 2018.

\section{Interferência e versões do tratamento}

O desenho supõe que o tratamento de um município não altera o desfecho de
outro e que o tratamento tem uma só versão. Quatro canais de interferência
são conhecidos: a deriva sobre municípios vizinhos e o aquífero compartilhado,
que atenuariam a estimativa; a migração do serviço aéreo para áreas sem
proibição, que a inflaria; e a mobilidade das mães. Como o análogo equatoriano não encontra efeito além de 150 metros do
bananal \cite{calzada2023bananas}, a interferência por deriva deve ser
pequena, mas os vizinhos imediatos dos tratados são separados dos demais,
formando três grupos. Quanto às versões, estar sob a lei significa coisas
diferentes: ter parado de voar; ter trocado o avião pela aplicação terrestre
ou pelo drone; ou não ter cumprido. O parâmetro, de intenção de tratar, é a
média sobre essas três versões. Não voar desde antes da lei não é uma quarta
versão: onde os voos cessaram antes de 2019, pela seca, pela regulação
municipal ou por decisão judicial, a lei não alterou a exposição, e $\beta$
atribuiria a ela o que tem outra causa.

\section{Análises secundárias e triangulação}
\label{sec:secundarias}

Três análises complementam o desenho principal. A primeira é a forma reduzida
estadual: um diferenças-em-diferenças com tratamento contínuo
\cite{callaway2024continuous}, em que a dose é a área plantada de banana antes
da lei \cite{ibgepam}, o grupo de dose nula é construído também pela aptidão
agroclimática e o alvo é o efeito no nível da dose recebida. Como a área mede
a cultura e não o método, classificar pela cultura produz falsos positivos, que
atenuam a comparação \cite{denteh2022misclassification}, e o Censo de 2006
sugere que há também falsos negativos; por isso a forma reduzida é lida como
descritiva. A segunda é o controle sintético e o
diferenças-em-diferenças sintético \cite{arkhangelsky2021sdid} para o agregado
dos municípios com avião, com doadores fora do alcance da interferência. A
terceira é a comparação com os municípios potiguares da Chapada do Apodi, que
não é experimento natural limpo, porque tudo o que é estadual muda junto na
divisa e as culturas diferem. A comparação entre as análises é descritiva, sem
rótulo de concordância: cada uma tem estimando próprio, e reporta-se o valor
que se esperaria de cada uma sob o efeito principal.

\section{Regra de decisão}
\label{sec:regra}

O Quadro~\ref{qua:regra} fixa o que cada resultado da medição e do cálculo de
poder aciona; o Apêndice~\ref{ap:regra} traz definições, critérios numéricos
e o enquadramento dos casos intermediários, entre eles os registros parciais
e os tardios. Os critérios foram escritos depois das rodadas exploratórias,
com os resultados delas à vista; nenhum é calculado com o desfecho do pós-lei. O efeito plausível é uma faixa: a
magnitude individual do análogo equatoriano, em três cenários que levam à
regra as ressalvas do Capítulo~\ref{cap:referencial}, multiplicada pela fração
de nascimentos exposta, em grade fixada no apêndice; decide a célula de 80
gramas e fração de 0,25.
Poder suficiente é condição necessária para o teste, não garantia de que o
efeito exista nessa magnitude. A data-limite para os registros é 31 de maio de
2027, fim do terceiro bimestre.

\begin{quadro}[htb]
\caption{\label{qua:regra}Regra de decisão}
\centering
\footnotesize
\begin{tabular}{@{}p{0.3cm}p{4.4cm}p{2.6cm}p{7.3cm}@{}}
\toprule
 & Medição dos voos & Poder & O que o ensaio entrega \\
\midrule
1 & Voos documentados por município e mês, com pico sazonal & Detectável $\leq$ plausível &
Desenho principal, confirmatório, com municípios e estação medidos. \\
2 & Idem & Detectável $>$ plausível &
Medida dos voos; desenho principal reportado como limite, sem teste. \\
3 & Voos documentados, sem pico sazonal & Qualquer &
Medida dos voos; o calendário não se aplica; forma reduzida descritiva. \\
4 & Série completa sem voos no período de referência & Qualquer &
Medida do alcance da lei; forma reduzida descritiva; sem estimativa de efeito da pulverização. \\
5 & Voos na referência e nenhum em 2018 (série completa): a prática cessou antes da lei & Qualquer &
Medida do alcance da lei; $\beta$ só descritivo, pois não distingue a lei do fim anterior da prática; sem estimativa de efeito da pulverização. \\
6 & Registros não obtidos, ou com menos de dois anos completos, até a data-limite & Com a classificação provisória &
Desenho principal com os municípios com aeronave no Censo de 2006 e a quadra chuvosa, exploratório (o parâmetro deste capítulo não é o alvo); cálculo de poder; medida parcial dos voos, se houver. \\
\bottomrule
\end{tabular}
\fonte{Elaboração própria. Definições e critérios no Apêndice~\ref{ap:regra}. A linha 4 prevalece sobre a 5, e a 5 sobre as linhas 1 a 3.}
\end{quadro}

\section{Inferência, poder e plano de análise}
\label{sec:inferencia}

Há poucos municípios tratados, contíguos, e a lei é atribuída no nível do
estado. Erros-padrão agrupados por município, com aproximação assintótica,
superestimam a precisão, e o \emph{wild cluster bootstrap} isolado não é
defensável com tão poucos tratados \cite{mackinnon2018wild}. O procedimento
principal é o de \citeonline{ferman2019inference}, que admite
heterocedasticidade por tamanho de município, aplicado à mudança do hiato
sazonal de cada município, líquida da chuva e dos efeitos fixos
(Apêndice~\ref{ap:regra}). Ele supõe municípios independentes, e os tratados
contíguos partilham choques; por isso o procedimento para um único agregado
tratado \cite{hagemann2025single} é reportado ao lado, e a conclusão
confirmatória exige que os dois concordem. Os de \citeonline{conley2011} e
\citeonline{canay2021wild} são sensibilidade, e a inferência conformal de
\citeonline{chernozhukov2021conformal} é usada no controle sintético.

Dois fatores limitam o poder e nenhum se resolve com mais anos de dados: o
número de municípios com aplicação aérea e a fração dos nascimentos de cada um
efetivamente exposta. O efeito mínimo detectável de $\beta$ tem nível de 5\%,
poder de 80\% e teste unilateral, sob o procedimento principal, e é calculado
apenas com dados do período de referência e do período de seca
(Apêndice~\ref{ap:regra}). Este projeto não traz o seu valor, o que é uma
limitação. A rodada exploratória de novembro, com classificação e estação
provisórias, indicou efeito detectável acima do plausível na célula que
decide, isto é, a linha 2 do Quadro~\ref{qua:regra}; a medição dos voos pode
mudar a classificação e o resultado. O valor pré-especificado é reportado no
plano de análise ou em adendo datado, antes da análise confirmatória.

\textbf{Rodadas exploratórias e plano de análise.} Duas rodadas exploratórias
antecedem este documento. Nenhuma estimativa delas é reportada; o
Apêndice~\ref{ap:regra} lista as variantes estimadas e o que estava fixado
antes de cada uma. A de setembro de 2026 estimou a forma reduzida estadual,
sob plano fechado no dia 21. A de novembro estimou o desenho de calendário,
ainda sem plano, com classificação e estação provisórias e sem a chuva, sobre
os nascimentos de 2005 a 2024: a exposição sazonal já foi cruzada com o
desfecho do pós-lei. Daí três compromissos. Municípios tratados e estação são
fixados pela medição dos voos, e os parâmetros ainda abertos, por critério
escrito sem uso do desfecho; a ratificação com o orientador não altera o
período de referência nem o Quadro~\ref{qua:regra}, e os limiares do
Apêndice~\ref{ap:regra} são provisórios até ela. Se municípios e estação medidos
coincidirem com variante já estimada, o resultado é rotulado como reestimação
em dados já examinados, e não como teste confirmatório. E o plano de análise
e o material das duas rodadas serão depositados em repositório público
versionado, com data verificável por terceiros, antes da análise
confirmatória, que é refeita a partir das bases originais.
